\chapter{Tokenization and sequence representation}

A model does not receive a DNA molecule, or even a string. It receives
numbers whose meaning comes from a representation contract. Before we
train a larger model, we must be able to explain every integer, every
tensor axis, and every base visible at a prediction step.

Our running sequence is \texttt{ACGTACG}. We will turn it into tokens,
integer IDs, and embedding vectors, then reconstruct it exactly. Along the
way we will build two negative controls: an overlapping-token objective
that partly reveals its own answer, and a strand transformation that breaks
when token boundaries are treated as biological boundaries.

The approach is implementation-first: predict a tiny result, derive it,
write the function, and test a counterexample. The programs and illustrations
are original. The author-maintained companion to Raschka's book informed
this read--implement--exercise progression, not our wording or figures
\cite{raschka}.

\section{Decide what a token means}

For base-level tokens, the sequence becomes seven tokens:
\texttt{A}, \texttt{C}, \texttt{G}, \texttt{T}, \texttt{A}, \texttt{C},
\texttt{G}. A \Term{k-mer} is a contiguous word of $k$ bases. For $k=3$,
we have at least two reasonable segmentations:
\begin{itemize}
\item overlapping windows: \texttt{ACG}, \texttt{CGT}, \texttt{GTA},
\texttt{TAC}, \texttt{ACG};
\item disjoint chunks with the short tail retained: \texttt{ACG},
\texttt{TAC}, \texttt{G}.
\end{itemize}
Both preserve the sequence if their decoding rules are known. They are
not interchangeable inputs to the same prediction objective.

\Plate{01-token-spans}{Three representations of \texttt{ACGTACG}. Colored bars denote source intervals, not chemical bonds. Overlapping windows share bases; disjoint chunks retain the terminal one-base token. Every token carries a half-open interval in the original sequence.}{fig:spans}

For a length-$L$ sequence with $L\geq k$, stride-one full windows give
\begin{equation}
T_{\mathrm{overlap}}=L-k+1.
\end{equation}
Disjoint chunks with a retained tail give
\begin{equation}
T_{\mathrm{disjoint}}=\left\lceil L/k\right\rceil.
\end{equation}
Our implementation additionally preserves short nonempty records with
$L<k$ as one short token; empty records produce none. The first formula
does not describe that fallback. Naming the edge case is better than forcing
a negative token count into a formula written for a different domain.

Each \texttt{Token} contains \texttt{text}, \texttt{start}, and
\texttt{stop}. Offsets count bases, not bytes or token positions.
\Listing{tokenize}

We accept uppercase canonical A/C/G/T only. Ambiguity symbols are rejected,
not silently mapped to a universal unknown token. Such a mapping would be
lossy: N, R, and a malformed character would become indistinguishable.
An ambiguity-aware extension should preserve the sets introduced in
DNA Computing Chapter 5 and explicitly define how the model consumes them.

\section{Build an integer code from first principles}

Assign digits $d(A)=0$, $d(C)=1$, $d(G)=2$, and $d(T)=3$. For a fixed-width
word $w=w_0\ldots w_{k-1}$, its base-four ID is
\begin{equation}
\mathrm{id}_k(w)=\sum_{j=0}^{k-1}d(w_j)4^{k-1-j}.
\label{eq:id}
\end{equation}
For \texttt{ACG}, this is $0\cdot16+1\cdot4+2=6$.
There are exactly $4^k$ fixed-width words and IDs $0$ through $4^k-1$.
The ID orders symbols; it does not measure biochemical similarity.

\Plate{02-base-four}{Encoding and decoding the same word. Horner's rule accumulates ID 6 from digits 0, 1, 2. Repeated division by four recovers remainders 2, 1, 0; reversing them restores ACG. Declared width preserves the leading A.}{fig:id}

\Listing{encode_kmer}
\Listing{decode_kmer}

The decoder needs $k$. ID zero could mean \texttt{A}, \texttt{AA}, or
\texttt{AAA} at different widths. A representation that discards width
cannot recover leading zero digits. Our tests enumerate every word for
$k=1,2,3,4$ and verify both directions of the bijection.

\subsection{A model vocabulary is not the raw k-mer code}

The teaching vocabulary reserves IDs 0, 1, and 2 for PAD, BOS, and EOS.
It then lists all one-base words, all two-base words, and so on through
width $k$, in base-four order. This policy retains every possible short tail.
For a width-$j$ word, the model ID is
\begin{equation}
\mathrm{model\_id}(w)
=3+\sum_{r=1}^{j-1}4^r+\mathrm{id}_j(w).
\label{eq:model-id}
\end{equation}
At $k=3$, \texttt{ACG} therefore has model ID $3+4+16+6=29$,
not 6. \texttt{TAC} has ID 72 and the short token \texttt{G} has ID 5.
Our disjoint record stores IDs $[29,72,5]$ alongside spans
$[0,3)$, $[3,6)$, and $[6,7)$.

Special symbols are reserved here; the sequence encoder does not insert
BOS or EOS, and the sequence decoder rejects all special IDs. A later
batching layer can add controls, but it must remove or interpret them
before calling this sequence-only decoder. Reserving an ID and actually
emitting it are different design decisions.

\section{Make the round trip enforce a contract}

For non-overlapping tokens, concatenation looks sufficient. For overlapping
tokens, it would duplicate bases: \texttt{ACG} followed by \texttt{CGT}
must reconstruct \texttt{ACGT}, not \texttt{ACGCGT}.
Offsets tell us exactly how much overlap should agree.

\Listing{reconstruct}

The function checks that each token extends a gap-free prefix, that its
declared length matches its text, and that overlapping bases agree.
For example, \texttt{ACG} at $[0,3)$ followed by \texttt{CTT} at $[1,4)$
must fail: the shared region is \texttt{CG} in the first token but
\texttt{CT} in the second. A decoder that simply takes the final character
would conceal the inconsistency.

This low-level function accepts any consistent extending spans. The higher
level \texttt{decode\_record} also checks that the decoded tokens equal the
declared overlap or disjoint segmentation of the reconstructed sequence.
It rejects mismatched vocabulary identities and control IDs. The record
therefore specifies not just a string, but how that string became model input.

We hash a canonical serialization of alphabet order, $k$, tail policy, and
the ordered vocabulary. The digest is an identity check, not authentication
against a malicious party. Store the mode and offsets in the record too.
Two vocabularies with equal size can give the same embedding row completely
different meanings; shape compatibility alone is not checkpoint compatibility.

\section{Count the representation cost}

With retained short words and three special symbols, vocabulary size is
\begin{equation}
V(k)=3+\sum_{j=1}^{k}4^j
=3+\frac{4^{k+1}-4}{3}.
\end{equation}
A float32 embedding table with width $D$ requires $4V(k)D$ bytes for its
parameter values. This excludes gradients, optimizer state, activations,
and allocation overhead.

\begin{center}\small\input{\Generated/resources.tex}\end{center}

The table uses $L=1000$ and $D=64$. At $k=6$, a disjoint representation
has 167 tokens instead of 1000, but the retained-tail vocabulary has 5,463
entries instead of seven at $k=1$. Its embedding values occupy 1,398,528
bytes. Overlapping width-six windows still have 995 tokens: a larger
vocabulary has barely shortened the sequence at all.

This is a parameter accounting exercise, not a throughput benchmark.
For an architecture with dense all-pairs attention, reducing token count
can reduce a quadratic interaction term; other costs scale differently.
A recurrent model, a convolution, and a state-space model have different
resource profiles. Tokenization should be evaluated against the actual
architecture and objective, not declared superior from one column.

\section{An embedding is a row lookup}

Let $E\in\mathbb{R}^{V\times D}$ be a learned table and let integer IDs
$I\in\{0,\ldots,V-1\}^{B\times T}$ index its rows. The output is
\begin{equation}
X_{b,t,:}=E_{I_{b,t},:},
\qquad X\in\mathbb{R}^{B\times T\times D}.
\label{eq:embedding}
\end{equation}
Equivalently, if $Q_{b,t,:}$ is one-hot over the vocabulary, the forward
calculation is $X_{b,t,:}=Q_{b,t,:}E$. We normally avoid materializing $Q$:
integer lookup expresses the operation without the mostly-zero tensor.
PyTorch documents the same input/output shape convention \cite{torch-embedding}.

\Plate{03-embedding}{A tiny lookup with repeated IDs. A is selected twice, so both positions receive the same row before positional information is added. Under a sum loss, its lookup gradients accumulate twice. PAD's assigned nonzero vector illustrates why a padding-gradient rule is not an attention or loss mask.}{fig:embedding}

For this isolated demonstration only, use the smaller ID scheme
A=0, C=1, G=2, T=3, PAD=4. It is deliberately separate from the retained-tail
vocabulary above. Set $E$ to the numbers 0 through 14 in a $5\times3$
matrix and look up $I=[0,2,0,4]$. The output is
\begin{equation}
X=
\begin{bmatrix}
0&1&2\\6&7&8\\0&1&2\\12&13&14
\end{bmatrix}.
\end{equation}
These assigned values make the calculation inspectable. They are not learned
biochemical coordinates or evidence that A and C differ by a distance of three.

\Listing{embedding_example}

For a scalar loss $\mathcal L=\sum X$, an ordinary row receives one unit
of gradient per occurrence in every feature:
\begin{equation}
\frac{\partial\mathcal L}{\partial E_{v,d}}
=\sum_{b,t:I_{b,t}=v}1,
\end{equation}
except for the explicit padding rule. The row gradients in our example
are $[2,2,2]$, $[0,0,0]$, $[1,1,1]$, $[0,0,0]$, and $[0,0,0]$.

The last row deserves close attention. PyTorch's \texttt{padding\_idx}
suppresses its gradient contribution from embedding lookup; it does not
force an explicitly supplied weight row to zero \cite{torch-embedding}.
Our PAD output remains $[12,13,14]$. Newly initialized
\texttt{nn.Embedding} padding rows default to zero, but that initialization
is not a universal guarantee about supplied weights or downstream operations.

The one-hot multiplication matches the \emph{forward values} here, not the
custom padding backward rule. Backpropagating through plain one-hot matrix
multiplication would give the selected PAD row a gradient too. A padding
ID also does not automatically mask attention, recurrent state updates, or
the training loss. Chapter 4's valid-target mask remains necessary.

\section{Track which bases a prediction can see}

Suppose the current input token is \texttt{ACG} and the next overlapping
token is \texttt{CGT}. Two of the three target bases, \texttt{CG}, are
already present in the input. This is not necessarily an invalid objective,
but its score must not be interpreted as predicting three unseen bases.

\Plate{04-target-reach}{The information boundary is a base coordinate. ACG ends at boundary 3. The next overlapping token CGT contains two already-observed bases and only one new base. Predicting the base at index 3 makes the unit of novelty explicit. A token-level causal mask cannot remove information already inside an input token.}{fig:reach}

Under an illustrative independent uniform base process, a fresh three-base
word has entropy $3\log 4=\log64$ nats. Conditional on the preceding
overlapping word, two bases are fixed; only the new base is uncertain, so
the conditional entropy is $\log4$. This is an analytical toy result,
not an empirical estimate of genomic entropy. Comparing those two losses
without the conditioning context would exaggerate the improvement.

For a transparent baseline, let each window predict the next unseen base.
Our function returns its input word, the target base, and the target index:
\Listing{safe_next_base_examples}

For the running sequence, the examples are
\texttt{ACG}$\to$\texttt{T} at 3,
\texttt{CGT}$\to$\texttt{A} at 4,
\texttt{GTA}$\to$\texttt{C} at 5, and
\texttt{TAC}$\to$\texttt{G} at 6.
The terminal \texttt{ACG} has no next-base target and is omitted.
This defines local training pairs; it does not yet implement a full
sequence model, its batching, or its evaluation split.

A useful negative control changes a future base while holding the prefix
fixed. The input representation for a prediction must remain unchanged.
Our test changes the base at index 3 and verifies that the first input token
is still \texttt{ACG} while its target changes. If a tokenizer were allowed
to include that base inside the current token, a later causal mask would
not repair the leak.

The same reasoning applies to masked-token objectives: masking one k-mer
may leave its bases visible inside neighboring overlapping k-mers. Define
the intended hidden base interval first, then remove every input feature
that reveals it. Also split raw records or biological groups before making
windows; random window splitting can put near-duplicates across evaluation
boundaries. Representation correctness is necessary but does not establish
evaluation independence.

\section{Reverse the strand, then reconsider the boundaries}

For full stride-one windows of width $k$ and $L\geq k$, reversing the token
order and reverse-complementing each token produces the windows of the
reverse-complemented sequence. A span $[a,b)$ becomes $[L-b,L-a)$,
as derived in DNA Computing Chapter 5.

\Plate{05-strand-boundaries}{Reverse complementation changes the sequence coordinate system. With overlapping full windows, reversing window order and each window's bases works. With disjoint chunks and a short tail, that operation leaves the short chunk at the front; re-tokenizing from the new origin gives different boundaries.}{fig:strand}

For our length-seven sequence, disjoint tokens are
\texttt{ACG|TAC|G}. Reversing token order and complementing gives
\texttt{C|GTA|CGT}. Concatenation correctly yields \texttt{CGTACGT},
but disjoint re-tokenization from the new origin yields
\texttt{CGT|ACG|T}. The sequence is right and the segmentation is wrong.
Token boundaries are a software choice, not covalent breaks in DNA.

Token index alone is also insufficient to identify a base position.
An overlapping width-three token at index $j$ starts at base $j$;
a disjoint full-width token at index $j$ starts at $3j$. Preserve base
spans and decide whether positions describe starts, centers, or something
else before adding positional features. A position-aware embedding can
distinguish equal tokens in different locations; the lookup table alone cannot.

Finally, consider \Term{canonicalization}: store the lexicographically
smaller of $s$ and $\operatorname{rc}(s)$. It merges \texttt{AAA} and
\texttt{TTT}. That may be suitable for a strand-invariant identity task,
but it loses orientation unless we keep a flip flag. Direction-dependent
labels must be transformed consistently. Even with a flag, a model that
never receives or uses it cannot learn two different labels for identical
canonicalized inputs. Palindromic ties use a deterministic no-flip convention.

\section{A representation checklist before training}

The completed artifact includes strict alphabet validation, fixed-width
integer codes, two segmentation policies, retained tails, offset-checked
reconstruction, vocabulary identity checks, an embedding gradient example,
and causal-reach tests. To reproduce the golden results and test the contracts:
\begin{Verbatim}[fontsize=\small]
python3 drafts/ch05/reference.py
python3 -m pytest tests/test_ch05_reference.py
\end{Verbatim}
The computed record is saved in \path{drafts/ch05/results.json}.

Before reusing the data with Chapter 6's recurrent models, answer four
questions. Can every record round-trip exactly? Does a checkpoint agree
on what each row means? Does each target lie beyond the allowed input
boundary? Are strand-dependent labels and coordinates still aligned?
Training cannot recover information a representation has discarded.

These examples do not train a genomic foundation model. DOGMA remains the
non-Transformer DNA-native research direction; Hermon DNA denotes the
Transformer-based direction; Evolutor is the broader framework. Tokenization
contracts apply to either architecture without proving either system exists
as a trained or validated biological model.

\Needspace{10\baselineskip}
\section{Exercises with worked solutions}

\subsection*{1. Predict the tail}
Tokenize \texttt{ACGTA} disjointly at $k=3$.
\textbf{Solution:} \texttt{ACG} at $[0,3)$ and \texttt{TA} at $[3,5)$.
Dropping the tail loses two bases; padding is a different policy requiring
its own length or mask metadata.

\subsection*{2. Encode without a table}
Find the raw width-three ID of \texttt{TGA}.
\textbf{Solution:} $3\cdot16+2\cdot4+0=56$. Its model ID in our retained-tail
vocabulary is $23+56=79$.

\subsection*{3. Recover leading zeros}
Decode ID 1 at widths one and three.
\textbf{Solution:} \texttt{C} and \texttt{AAC}. The integer alone does not
encode width.

\subsection*{4. Detect inconsistent overlap}
Why reject \texttt{ACG} at $[0,3)$ followed by \texttt{CTT} at $[1,4)$?
\textbf{Solution:} shared bases should be \texttt{CG} but the second token
claims \texttt{CT}. Taking only the new suffix would hide the disagreement.

\subsection*{5. Count vocabulary entries}
Compute the retained-tail vocabulary for $k=2$.
\textbf{Solution:} $3+4+16=23$. At embedding width 64 in float32, parameter
values occupy $23\cdot64\cdot4=5{,}888$ bytes, excluding training state.

\subsection*{6. Follow a repeated lookup}
For IDs $[2,2,2]$ and a sum loss, what gradient reaches row 2?
\textbf{Solution:} three in every feature, unless that row is the declared
padding index. Repeated lookup shares parameters; it does not create three
independent rows.

\subsection*{7. Separate padding mechanisms}
Does a zero-gradient padding row imply padding cannot affect a model?
\textbf{Solution:} no. Its forward vector can be nonzero, positional features
can alter it, and downstream operations still need appropriate masks.
Lookup-gradient suppression is not a universal padding exclusion rule.

\subsection*{8. Measure novelty}
How many unseen bases are in the next stride-one width-six token?
\textbf{Solution:} one; five overlap. Under independent uniform bases,
its conditional uncertainty is $\log4$, not $\log4096$ nats.

\subsection*{9. Audit a causal mask}
A token covering $[0,3)$ is used to predict base 2. Can a token-level
causal mask fix the problem? \textbf{Solution:} no. Base 2 is already inside
the token. Change token reach or move the target beyond the token's stop.

\subsection*{10. Reverse an uneven partition}
Reverse-complement the tokens \texttt{ACG|TAC|G}, then re-tokenize.
\textbf{Solution:} transformed chunks \texttt{C|GTA|CGT} concatenate
correctly, but the new disjoint partition is \texttt{CGT|ACG|T}.

\subsection*{11. Preserve directional labels}
A dataset labels \texttt{AAA} and \texttt{TTT} differently. What happens
after canonicalization without orientation metadata?
\textbf{Solution:} both become \texttt{AAA}. A deterministic predictor sees
identical inputs with incompatible labels. Keep and use orientation, transform
labels appropriately, or choose a genuinely strand-invariant task.

\subsection*{12. Specify a checkpoint contract}
What must accompany embedding weights?
\textbf{Worked rubric:} vocabulary ordering and identity, alphabet policy,
special-ID meaning, segmentation and tail policy, positional convention,
architecture dimensions, and preprocessing version. Shape checks are necessary
but insufficient. Add a golden raw-sequence-to-ID example and reject a
deliberately permuted vocabulary even when its dimensions match.
